\section{Related Work}\label{sec:related}
\textbf{Federated IDS.} A survey~\cite{agrawal2022} documents the uptake of FL-IDS; FLAD proposes adaptive federated learning for DDoS detection~\cite{doriguzzi2024}. Devine et al.\ train a federated SVM on IoT DDoS data with a dedicated data-preparation stage~\cite{devine2025}.
\textbf{Non-IID clients.} Label-skew partitions are commonly drawn from Dirichlet priors~\cite{hsu2019}; FedProx adds a proximal term against client drift~\cite{li2020fedprox}. FL-IDS variants use knowledge distillation with a proximal term~\cite{fdids2025}, dynamic weighted clustering~\cite{fedDWC2025} and multi-dataset benchmarks of FedAvg/FedProx/FedNova~\cite{datasetcentric2026}; most change the optimiser or model, not the preprocessing.
\textbf{Selection and statistics sharing.} Filter selection for IDS is classical~\cite{ambusaidi2016}; federated variants use greedy per-attack selection~\cite{yang2021}, chi-square with a GRU~\cite{nguyen2026} or mutual information~\cite{fdids2025}. StatAvg~\cite{bouzinis2025} shares per-feature means and variances so that all clients normalise identically, reporting accuracy gains of 4--17\% over normalisation-based baselines such as FedBN. We keep its principle but share per-class second moments, which suffice for Pearson and LDA.
\textbf{Evaluation pitfalls.} Preprocessing fitted before splitting and duplicates are recurrent sources of optimism in ML-based security~\cite{arp2022,kapoor2023}. Table~\ref{tab:related} (also~\ref{tab:lit}) compares the closest methods and quotes their results (not directly comparable).
\begin{table*}[t]
\centering
\caption{Related-work comparison and reported results (as far as the cited abstracts/papers state; ``part'' = partial or not stated). Reported numbers are quoted from each source; datasets, splits and metrics differ, so rows are not directly comparable.}
\label{tab:related}\label{tab:lit}
\tiny
\setlength{\tabcolsep}{2.4pt}
\begin{tabular}{p{2.3cm}p{1.9cm}p{1.6cm}p{1.6cm}p{0.6cm}p{0.8cm}p{1.0cm}p{1.1cm}p{3.6cm}}
\toprule
Method & Dataset(s) & Model & Selection & Fed. & Non-IID & Fed. prep. stats & Leakage-free audit & Reported result (quoted) \\
\midrule
Karunamurthy et al. \cite{karunamurthy2025} & MQTTset & CNN & wrapper (COA) & yes & no & no & no & acc. 95.59\%, macro-F1 94.04\% \\
Devine et al. \cite{devine2025} & IoT DDoS data & SVM/RF/ANN & data preparation & yes & no & no & part & not stated in abstract \\
Yang \& Kondo \cite{yang2021} & NSL-KDD & on-device NN & greedy per attack & yes & no & no & no & not stated in abstract \\
FD-IDS \cite{fdids2025} & Edge-IIoT, N-BaIoT & NN + KD + prox & mutual information & yes & yes & no & no & not stated in abstract \\
StatAvg \cite{bouzinis2025} & TON-IoT, CIC-IoT-2023 & deep FL & -- (normalisation) & yes & yes & part & no & acc. gain 4--17\% over normalisation baselines \\
Fed-FeRe \cite{nguyen2026} & IoT IDS data & GRU & chi-square & yes & yes & no & no & +3.06\% precision, $\approx$2\% less traffic than FedAvg \\
FedDWC \cite{fedDWC2025} & IoTID20 & clustered FL & -- & yes & yes & no & no & acc. +1.9/+1.31/+1.01\% over FedAvg/FedProx/IFCA \\
Senior \cite{senior2026} & CIC-IoT2023 & linear SVM & Pearson$\cap$LDA & yes & no & no & no & F1 98.81\% (99.66\% in our re-run) \\
FedStat-FS (ours) & CIC-IoT2023, N-BaIoT & linear SVM (+MLP/prox baselines) & Pearson$\cap$LDA, exact & yes & yes & yes & yes & Table~\ref{tab:main} \\
\bottomrule
\end{tabular}
\end{table*}
